\subsection{Сила Ампера}

%%%%%%%%%%%%%%%%%%%%%%%%%%%%%%%%%%%%%%%%%%%%%%%%%%%%
%%%%%%%%%%%%%%%%%%%%%%%%%%%%%%%%%%%%%%%%%%%%%%%%%%%%
 \z{На горизонтальных рельсах, расположенных в вертикальном однородном магнитном поле $B$, покоится металлическая перемычка перпендикулярно рельсам. Расстояние между рельсами $l$, масса перемычки $m$, коэффициент трения между рельсами и перемычкой равен $\mu$. Какой ток нужно пропустить через перемычку, чтобы она сдвинулась с места?}
%
{$I = \frac{\mu m g}{B l}$}

%%%%%%%%%%%%%%%%%%%%%%%%%%%%%%%%%%%%%%%%%%%%%%%%%%%%
 \z{Проволочный квадратик с током находится в однородном магнитном поле, индукция которого $B = 0.01\,\text{Тл}$. Сила тока в кольце $I = 0.5\,\text{А}$. Сторона квадратика $a = 2\,\text{см}$. Какой максимальный момент сил может действовать на кольцо со стороны магнитного поля?}
%
{$M = I B a^{2} = 2 \cdot 10^{-6}\,\text{Н}\cdot\text{м}$}

%%%%%%%%%%%%%%%%%%%%%%%%%%%%%%%%%%%%%%%%%%%%%%%%%%%%
 \z{Найдите максимальный момент сил, который может действовать на плоский контур площадью $S$, находящийся в магнитном поле $B$, в котором течет ток $I$.}
%
{$M = I B S$}

%%%%%%%%%%%%%%%%%%%%%%%%%%%%%%%%%%%%%%%%%%%%%%%%%%%%
 \z[{\begin{wrapfigure}[4]{r}{0.15\textwidth}%
		\vspace{-25pt}%
		\begin{flushright}%
 			\includegraphics[width=0.15\textwidth]{5.2.4}%
		\end{flushright}%
	\end{wrapfigure}%
	}]%
{На линейный проводник длины $l$, расположенный перпендикулярно магнитному полю, действует сила $F$, если ток в проводнике равен $I$. С какой силой магнитное поле будет действовать на изогнутый под углом $30\degree$ проводник длины $l + L$, если плоскость изгиба перпендикулярна магнитному полю, а ток в проводнике равен $2I$?}
%
{$F = 2F \sqrt{1 + \frac{L^{2}}{l^{2}} - 2 \frac{L}{l} \cos\varphi}$}

%%%%%%%%%%%%%%%%%%%%%%%%%%%%%%%%%%%%%%%%%%%%%%%%%%%%
 \z[{\begin{wrapfigure}[5]{r}{0.15\textwidth}%
		\vspace{-15pt}%
		\begin{flushright}%
 			\includegraphics[width=0.15\textwidth]{5.2.5}%
		\end{flushright}%
	\end{wrapfigure}%
	}]%
{Кольцо радиуса $R$, по которому течёт ток $I$, поместили в неоднородное аксиально-симметричное поле. Ось кольца совпадает с осью симметрии магнитного поля. Индукция магнитного поля $B$, действующего на ток, направлена под углом $30\degree$ к оси симметрии поля. Масса кольца $m$. Определите ускорение кольца.}
%
{$a = \frac{2\pi R B I}{m} \sin\alpha$}

%%%%%%%%%%%%%%%%%%%%%%%%%%%%%%%%%%%%%%%%%%%%%%%%%%%%
 \z[{\begin{wrapfigure}[6]{r}{0.40\textwidth}%
		\vspace{-15pt}%
		\begin{flushright}%
 			\includegraphics[width=0.40\textwidth]{5.2.6}%
		\end{flushright}%
	\end{wrapfigure}%
	}]%
{На непроводящей горизонтальной поверхности стола лежит проводящая жёсткая тонкая квадратная рамка из однородного куска проволоки со стороной, равной $a$. Рамка находится в магнитном поле длинного горизонтального провода с током, расположенного симметрично над рамкой. Масса рамки $m$, индукция магнитного поля у боковых сторон рамки 1 и 2 равна $B$. Коэффициент трения скольжения рамки о поверхность стола равен $\mu$ ($\mu < \frac{1}{3}$). Какой силы ток нужно пропустить по рамке, чтобы она начала скользить по столу, не отрываясь от него?}
%
{$I \ge \frac{\mu m g}{a B}$}

%%%%%%%%%%%%%%%%%%%%%%%%%%%%%%%%%%%%%%%%%%%%%%%%%%%%
 \z[{\begin{wrapfigure}[5]{r}{0.15\textwidth}%
		\vspace{-15pt}%
		\begin{flushright}%
 			\includegraphics[width=0.15\textwidth]{5.2.7}%
		\end{flushright}%
	\end{wrapfigure}%
	}]%
{На непроводящей горизонтальной поверхности стола лежит проводящая жёсткая тонкая рамка из однородного куска проволоки в виде равностороннего треугольника со стороной, равной $a$. Рамка находится в однородном горизонтальном магнитном поле, линии индукции которого перпендикулярны одной из сторон рамки. Масса рамки $m$, величина индукции $B$. Какой силы ток нужно пропустить по рамке, чтобы она начала приподниматься относительно одной из вершин треугольника?}
%
{$I = \frac{4 m g}{3 a B}$}






